\documentclass[../../main.tex]{subfiles}
\begin{document}
\section{习题}\label{section:ch14-problems}

\begin{exercise}\label{exercise:ch14-01}
证明：有限维向量空间上的余向量 $\omega^1,\dots,\omega^k$ 线性相关当且仅当 $\omega^1\wedge\cdots\wedge\omega^k=0$。
\end{exercise}

\begin{exercise}\label{exercise:ch14-02}
对哪些 $k$ 与 $n$，$\mathbb{R}^n$ 上每个 $k$-余向量都可分解？[建议：先做 $n=3$ 与 $n=4$ 的情形，然后看你能否把结果推广到其他维数。]
\end{exercise}

\begin{exercise}\label{exercise:ch14-03}
我们有两种方式来思考有限维向量空间 $V$ 上的协变 $k$-张量：具体地说，把它们视为 $V$ 上的 $k$-多重线性泛函；抽象地说，把它们视为抽象张量积空间 $V^*\otimes\cdots\otimes V^*$ 的元素。然而我们只借助具体定义定义了斜对称张量与对称张量。本题概述一种处理斜对称张量的抽象途径。（对称张量可类似处理。）

设 $V$ 是一个有限维实向量空间。令 $A$ 表示 $k$ 重抽象张量积 $V^*\otimes\cdots\otimes V^*$ 中由所有形如 $\omega^1\otimes\cdots\otimes\omega^k$（其中对某对 $i\ne j$ 有 $\omega^i=\omega^j$）的元素张成的子空间。（因此当 $k<2$ 时 $A$ 是平凡子空间。）令 $\Lambda^k(V^*)$ 表示商向量空间 $(V^*\otimes\cdots\otimes V^*)/A$。
\begin{enumerate}[(a)]
\item 证明存在唯一同构 $F:\Lambda^k(V^*)\to\Lambda^k(V^*)$，使下图交换：
\begin{align*}
\begin{array}{ccc}
V^*\otimes\cdots\otimes V^*&\xrightarrow{\cong}&T^k(V^*)\\
\downarrow\pi&&\downarrow\operatorname{Alt}\\
\Lambda^k(V^*)&\xrightarrow{F}&\Lambda^k(V^*)
\end{array}
\end{align*}
其中 $\pi:V^*\otimes\cdots\otimes V^*\to\Lambda^k(V^*)$ 是投影。
\item 在 $\bigoplus_k\Lambda^k(V^*)$ 上定义外积 $\omega\wedge\eta=\pi(\widetilde\omega\otimes\widetilde\eta)$，其中 $\widetilde\omega,\widetilde\eta$ 是满足 $\pi(\widetilde\omega)=\omega$、$\pi(\widetilde\eta)=\eta$ 的任意张量。证明此外积良定义，且 $F$ 把此外积映为 $\Lambda(V^*)$ 上 Alt 约定的外积。
\end{enumerate}
[注：这就是为什么有些作者认为外积的 Alt 约定比行列式约定更自然的原因之一。]
\end{exercise}

\begin{exercise}\label{exercise:ch14-04b}
本章由于斜对称协变张量在微分几何中有许多重要应用而着重讨论它们。斜对称反变张量也有一些应用（习题\ref{exercise:ch21-14}描述了其中一个）。若 $V$ 是一个有限维向量空间，我们可以把斜对称反变 $k$-张量定义为 $V^*\times\cdots\times V^*\to\mathbb{R}$ 中交换任意两个自变量就变号的多重线性泛函，或定义为某个商空间 $(V\otimes\cdots\otimes V)/A$（类似于习题\ref{exercise:ch14-03}中所定义者）的元素。在第一种情形下，外积的定义恰如\eqref{align:ch14-wedge-def}，只是把向量与余向量的角色互换。在第二种情形下，外积定义为张量积在商空间中的像，如习题\ref{exercise:ch14-03}。无论采用哪种定义，斜对称反变 $k$-张量都称为多重向量或 $k$-向量。对本题，选择你偏好的任一一种定义。
\begin{enumerate}[(a)]
\item 证明 $V$ 中元素的有序 $k$ 元组 $(v_1,\dots,v_k)$ 线性相关当且仅当 $v_1\wedge\cdots\wedge v_k=0$。
\item 证明两个线性无关的有序 $k$ 元组 $(v_1,\dots,v_k)$ 与 $(w_1,\dots,w_k)$ 张成相同子空间，当且仅当存在某个非零实数 $c$，使
\begin{align*}
v_1\wedge\cdots\wedge v_k=c\,w_1\wedge\cdots\wedge w_k.
\end{align*}
\end{enumerate}
\end{exercise}

\begin{exercise}[Cartan 引理]\label{exercise:ch14-05}
设 $M$ 是一个带边或不带边的光滑 $n$ 流形，$(\omega^1,\dots,\omega^k)$ 是开子集 $U\subseteq M$ 上一个有序光滑 $1$-形式组，使得对每个 $p\in U$，$(\omega^1|_p,\dots,\omega^k|_p)$ 都线性无关。给定 $U$ 上光滑 $1$-形式 $\alpha^1,\dots,\alpha^k$，满足
\begin{align*}
\sum_{i=1}^k\alpha^i\wedge\omega^i=0,
\end{align*}
证明每个 $\alpha^i$ 都可写成 $\omega^1,\dots,\omega^k$ 带光滑系数的线性组合。
\end{exercise}

\begin{exercise}\label{exercise:ch14-06}
由
\begin{align*}
\omega=x\,\mathrm{d}y\wedge\mathrm{d}z+y\,\mathrm{d}z\wedge\mathrm{d}x+z\,\mathrm{d}x\wedge\mathrm{d}y
\end{align*}
在 $\mathbb{R}^3$ 上定义一个 $2$-形式。
\begin{enumerate}[(a)]
\item 在由 $(x,y,z)=(\rho\sin\varphi\cos\theta,\rho\sin\varphi\sin\theta,\rho\cos\varphi)$ 定义的球坐标 $(\rho,\varphi,\theta)$ 下计算 $\omega$。
\item 在直角坐标与球坐标下分别计算 $\mathrm{d}\omega$，并验证两个表达式表示同一个 $3$-形式。
\item 利用 $(\varphi,\theta)$ 在该坐标有定义的开子集上的坐标，计算 $S^2$ 的包含映射 $\iota:S^2\hookrightarrow\mathbb{R}^3$ 的拉回 $\iota^*\omega$。
\item 证明 $\iota^*\omega$ 处处不为零。
\end{enumerate}
\end{exercise}

\begin{exercise}\label{exercise:ch14-07}
在下列各情形中，$M,N$ 都是光滑流形；$F:M\to N$ 是一个光滑映射；$\omega$ 是 $N$ 上的一个光滑微分形式。在每种情形中，计算 $\mathrm{d}\omega$ 与 $F^*\omega$，并通过直接计算验证 $F^*(\mathrm{d}\omega)=\mathrm{d}(F^*\omega)$。（参见习题\ref{exercise:ch11-07}。）
\begin{enumerate}[(a)]
\item $M=N=\mathbb{R}^2$；$F(s,t)=(st,\mathrm{e}^t)$；$\omega=x\,\mathrm{d}y$。
\item $M=\mathbb{R}^2$，$N=\mathbb{R}^3$；$F(\theta,\varphi)=((\cos\varphi+2)\cos\theta,(\cos\varphi+2)\sin\theta,\sin\varphi)$；$\omega=y\,\mathrm{d}z\wedge\mathrm{d}x$。
\item $M=\{(u,v)\in\mathbb{R}^2:u^2+v^2<1\}$，$N=\mathbb{R}^3\setminus\{0\}$；$F(u,v)=\left(u,v,\sqrt{1-u^2-v^2}\right)$；$\omega=\dfrac{x\,\mathrm{d}y\wedge\mathrm{d}z+y\,\mathrm{d}z\wedge\mathrm{d}x+z\,\mathrm{d}x\wedge\mathrm{d}y}{(x^2+y^2+z^2)^{3/2}}$。
\end{enumerate}
\end{exercise}

\begin{exercise}\label{exercise:ch14-08}
对每个非负整数 $k$，证明存在一个反变函子 $\Omega^k:\mathbf{Diff}\to\mathbf{Vec}_{\mathbb R}$，它把每个光滑流形 $M$ 映到向量空间 $\Omega^k(M)$，把每个光滑映射 $F$ 映到拉回 $F^*$。证明外导数是从 $\Omega^k$ 到 $\Omega^{k+1}$ 的一个自然变换。（参见习题\ref{exercise:ch11-18}。）
\end{exercise}

\begin{exercise}\label{exercise:ch14-09}
设 $M,N$ 是光滑流形，$\pi:M\to N$ 是一个满射光滑淹没，且纤维连通。若 $v\in T_pM$ 满足 $\mathrm{d}\pi_p(v)=0$，就称 $v$ 是\emph{竖直的}。设 $\omega\in\Omega^k(M)$。证明：存在 $\eta\in\Omega^k(N)$ 使 $\omega=\pi^*\eta$，当且仅当对每个 $p\in M$、每个竖直向量 $v\in T_pM$，都有 $v\lrcorner\omega_p=0$ 且 $v\lrcorner\mathrm{d}\omega_p=0$。[提示：先做 $\pi:\mathbb{R}^{n+m}\to\mathbb{R}^n$ 为到前 $n$ 个坐标的投影这一情形。]
\end{exercise}
\end{document}
